% !TeX program = xelatex
\documentclass[aspectratio=169,10pt]{beamer}
\usepackage{amsmath,amssymb,array,booktabs,mathtools}
\usepackage{graphicx,hyperref}
\usepackage{tikz}
% ============================================================
% Beamer Color System — Speculative Decoding black theme
% ============================================================

\providecommand{\beamerthemename}{black}

\RequirePackage{xifthen}

% Keep the public theme selector compatible with the source deck. Its current
% variants intentionally resolve to the same monochrome academic palette.
\newcommand{\usemonochrometheme}{%
  \definecolor{themePrimary}{HTML}{111111}%
  \definecolor{themeSecondary}{HTML}{444444}%
  \definecolor{themeTertiary}{HTML}{000000}%
  \definecolor{themeLight}{HTML}{B8B8B8}%
  \definecolor{themeSilver}{HTML}{E6E6E6}%
  \definecolor{themeBg}{HTML}{FAFAFA}%
  \definecolor{themeBorder}{HTML}{D7D7D7}%
  \definecolor{themeAlert}{HTML}{111111}%
  \definecolor{themeExample}{HTML}{111111}%
}

\ifthenelse{\equal{\beamerthemename}{black}}{\usemonochrometheme}{%
\ifthenelse{\equal{\beamerthemename}{blue}}{\usemonochrometheme}{%
\ifthenelse{\equal{\beamerthemename}{gold}}{\usemonochrometheme}{%
\ifthenelse{\equal{\beamerthemename}{green}}{\usemonochrometheme}{%
\ifthenelse{\equal{\beamerthemename}{purple}}{\usemonochrometheme}{%
  \PackageError{beamer-colors}{Unknown theme '\beamerthemename'.%
    Valid themes: black, blue, gold, green, purple}{}
}}}}}

\setbeamercolor{structure}{fg=themePrimary}
\setbeamercolor{palette primary}{bg=themePrimary,fg=white}
\setbeamercolor{palette secondary}{bg=themeSecondary,fg=white}
\setbeamercolor{palette tertiary}{bg=themeTertiary,fg=white}
\setbeamercolor{palette quaternary}{bg=themePrimary!80!black,fg=white}

\setbeamercolor{title}{fg=themePrimary}
\setbeamercolor{frametitle}{fg=themePrimary}
\setbeamercolor{subtitle}{fg=themeSecondary}
\setbeamercolor{author}{fg=themePrimary!80!black}
\setbeamercolor{institute}{fg=themeSecondary}
\setbeamercolor{date}{fg=themeSecondary}

\setbeamercolor{normal text}{fg=black!84,bg=white}
\setbeamercolor{alerted text}{fg=themeAlert}
\setbeamercolor{example text}{fg=themeExample}

\setbeamercolor{block title}{bg=themePrimary,fg=white}
\setbeamercolor{block body}{bg=themeSilver!40}
\setbeamercolor{block title alerted}{bg=themeAlert,fg=white}
\setbeamercolor{block body alerted}{bg=themeAlert!8}
\setbeamercolor{block title example}{bg=themeExample,fg=white}
\setbeamercolor{block body example}{bg=themeExample!8}

\setbeamercolor{itemize item}{fg=themePrimary}
\setbeamercolor{itemize subitem}{fg=themeSecondary}
\setbeamercolor{itemize subsubitem}{fg=themeSecondary!70}
\setbeamercolor{enumerate item}{fg=themePrimary}
\setbeamercolor{enumerate subitem}{fg=themeSecondary}

\setbeamercolor{section in toc}{fg=themePrimary}
\setbeamercolor{subsection in toc}{fg=themeSecondary}
\setbeamercolor{footnote mark}{fg=themePrimary}
\setbeamercolor{caption name}{fg=themePrimary}

\definecolor{positive}{HTML}{111111}
\definecolor{negative}{gray}{0.42}
\definecolor{emphasis}{HTML}{111111}
\definecolor{neutral}{gray}{0.55}

\newcommand{\pos}[1]{\textcolor{positive}{#1}}
\newcommand{\con}[1]{\textcolor{negative}{#1}}
\newcommand{\HL}[1]{\textcolor{emphasis}{#1}}

\makeatletter
\@ifpackageloaded{hyperref}{%
  \hypersetup{
    colorlinks=true,
    linkcolor=black!75,
    urlcolor=black!75,
    citecolor=black!75,
  }%
}{}
\makeatother

% ============================================================
% Layout: Metropolis — Modern minimal, no external dependencies
% Recreates the metropolis aesthetic using standard beamer.
% ============================================================

\usetheme{default}
\usecolortheme{default}

\newlength{\bsProgressDone}
\newlength{\bsProgressRemain}

% ---- Typography ----
\usefonttheme{professionalfonts}
\setbeamerfont{title}{size=\Large,series=\bfseries}
\setbeamerfont{frametitle}{size=\large,series=\bfseries}
\setbeamerfont{framesubtitle}{size=\normalsize,series=\mdseries}
\setbeamerfont{author}{size=\normalsize}
\setbeamerfont{institute}{size=\small}
\setbeamerfont{date}{size=\small}
\setbeamerfont{section title}{size=\Large,series=\bfseries}
\setbeamerfont{footline}{size=\tiny}

% ---- Remove navigation ----
\setbeamertemplate{navigation symbols}{}

% ---- Frame title: left-aligned, clean ----
\setbeamertemplate{frametitle}{%
  \vskip6pt%
  \usebeamerfont{frametitle}\usebeamercolor[fg]{frametitle}\insertframetitle%
  \ifx\insertframesubtitle\empty\else
    \\[2pt]\usebeamerfont{framesubtitle}\usebeamercolor[fg]{subtitle}\insertframesubtitle%
  \fi
  \vskip4pt%
  {\color{themePrimary}\hrule height 0.45pt}%
  \vskip6pt%
}

% ---- Footline: thin progress bar + page number ----
\setbeamertemplate{footline}{%
  \nointerlineskip%
  \pgfmathsetmacro{\progressfrac}{\insertframenumber/\inserttotalframenumber}%
  \pgfmathsetlength{\bsProgressDone}{\progressfrac\paperwidth}%
  \pgfmathsetlength{\bsProgressRemain}{\paperwidth-\bsProgressDone}%
  \hbox to\paperwidth{%
    {\color{themePrimary!75}\rule{\bsProgressDone}{1pt}}%
    {\color{themeBorder!65}\rule{\bsProgressRemain}{1pt}}%
  }%
  \vskip2pt%
  \hbox to\paperwidth{\hfill\usebeamerfont{footline}\textcolor{black!45}{\insertframenumber\,/\,\inserttotalframenumber}\hspace{8pt}}%
  \vskip4pt%
}

% ---- Title page: minimal ----
\setbeamertemplate{title page}{%
  \vfill
  \begin{flushleft}
    {\usebeamerfont{title}\usebeamercolor[fg]{title}\inserttitle}\\[18pt]
    {\usebeamerfont{author}\textcolor{black!60}{\insertauthor}}\\[4pt]
    {\usebeamerfont{date}\textcolor{black!60}{\insertdate}}
  \end{flushleft}
  \vfill
}

% ---- Section page ----
\AtBeginSection{%
  \begin{frame}[plain,noframenumbering]
    \vfill
    \begin{center}
      {\usebeamerfont{section title}\usebeamercolor[fg]{structure}\insertsectionhead}
    \end{center}
    \vfill
  \end{frame}
}

% ---- Itemize: clean circles ----
\setbeamertemplate{itemize item}{\small\raise1pt\hbox{\textbullet}}
\setbeamertemplate{itemize subitem}{\scriptsize\raise1pt\hbox{--}}
\setbeamertemplate{itemize subsubitem}{\scriptsize\raise1pt\hbox{$\cdot$}}

\author{Skill QA}
\date{5 September 2026}
\newcommand{\slidecite}[1]{%
  \begin{tikzpicture}[remember picture,overlay]
    \node[anchor=south west,align=left,text width=0.92\paperwidth,
      inner sep=0pt,font=\tiny,text=black!38]
      at ([xshift=0.68cm,yshift=0.72cm]current page.south west)
      {\raggedright #1\par};
  \end{tikzpicture}%
}
\newcommand{\norm}[1]{\left\lVert #1\right\rVert}

\title{Waiting Near Capacity}
\hypersetup{pdftitle={Waiting Near Capacity},pdfauthor={Skill QA}}
\begin{document}
\begin{frame}[plain]
  \titlepage
\end{frame}

\begin{frame}{Random arrivals share one server}
  \vspace{0.15cm}
  \begin{center}
    {\large
    \(\underbrace{\text{arrivals}}_{\lambda=6\ \text{jobs/s}}
    \ \longrightarrow\ \text{waiting line}
    \ \longrightarrow\
    \underbrace{\text{one server}}_{\mu=10\ \text{jobs/s}}
    \ \longrightarrow\ \text{departures}\)}
  \end{center}
  \vspace{0.35cm}
  \begin{itemize}
    \item Poisson arrivals and independent exponential service times.
    \item First-come first-served, unbounded waiting room.
    \item Stationary regime: \(\lambda<\mu\).
  \end{itemize}
  \vspace{0.25cm}
  \(\lambda\): average arrival rate.\qquad \(\mu\): service rate while busy.
  \slidecite{Supplied toy fixture Q1; explicit model assumptions A1 in source-notes.md.}
\end{frame}

\begin{frame}{System time includes 0.15 seconds of waiting}
  \vspace{0.05cm}
  \[
    \underbrace{\mathbb{E}[T]}_{\text{arrival to departure}}
    =\frac{1}{\mu-\lambda}
    =\frac{1}{10-6}
    =\boxed{0.25\ \text{s}}
  \]
  \vspace{0.35cm}
  \begin{center}
    \renewcommand{\arraystretch}{1.5}
    \begin{tabular}{lcc}
      \toprule
      Mean time & Calculation & Seconds\\
      \midrule
      Service & \(1/\mu\) & \(0.10\)\\
      Waiting & \(0.25-0.10\) & \(\mathbf{0.15}\)\\
      \bottomrule
    \end{tabular}
  \end{center}
  \vspace{0.25cm}
  \begin{center}
    Utilization is \(\rho=\lambda/\mu=0.60\), yet arrivals can still wait.
  \end{center}
  \slidecite{Supplied formula Q1; calculation D1 in source-notes.md. Model means, not measured latencies.}
\end{frame}

\begin{frame}{Less spare capacity produces sharply higher delay}
  \vspace{0.1cm}
  \begin{center}
    \(\mu=10\) jobs/s throughout.\qquad Mean service time stays at \(0.10\) s.
    \vspace{0.4cm}

    \renewcommand{\arraystretch}{1.5}
    \begin{tabular}{rccc}
      \toprule
      \(\lambda\) (jobs/s) & Utilization & System time (s) & Waiting (s)\\
      \midrule
      2 & \(20\%\) & \(0.125\) & \(0.025\)\\
      \textbf{6} & \(\mathbf{60\%}\) & \(\mathbf{0.250}\) & \(\mathbf{0.150}\)\\
      8 & \(80\%\) & \(0.500\) & \(0.400\)\\
      9 & \(90\%\) & \(1.000\) & \(0.900\)\\
      \bottomrule
    \end{tabular}
  \end{center}
  \vspace{0.2cm}
  \begin{center}
    From 6 to 9 jobs/s, mean waiting grows sixfold.
  \end{center}
  \slidecite{Calculated sensitivity example D2 in source-notes.md. Every row assumes stationary M/M/1.}
\end{frame}

\begin{frame}{Near capacity, the model's mean waiting diverges}
  \vspace{0.1cm}
  \[
    \underbrace{\mathbb{E}[W_q]}_{\text{mean waiting}}
    =\frac{1}{\mu-\lambda}-\frac{1}{\mu}
    =\frac{\lambda}{\mu(\mu-\lambda)}
    \ \longrightarrow\ \infty
    \quad\text{as }\lambda\to\mu^{-}.
  \]
  \vspace{0.45cm}
  \begin{itemize}
    \item Random arrivals and service create backlogs. Less spare capacity leaves less room to clear them.
    \item High utilization can coexist with large waiting times.
  \end{itemize}
  \vspace{0.35cm}
  {\color{black!55}Scope: stationary means under this model. No tail-latency or deployment guarantee.}
  \slidecite{Algebra and limit D3; assumptions A1 in source-notes.md. The finite-mean formula requires \(\lambda<\mu\).}
\end{frame}
\end{document}
